% ============================================================
%  Section: Chemical Organization Theory and the
%           Statistical Geometry of Regimes
%  Label:   sec:cot_lp
%  Figures: hierarchy_complex_orgs.png       (from HTML export)
%  Inter-org (outputs/complex_COT_ESMO/):
%           fig_inter1_heatmaps.png           (cover-edge 2x2 matrices)
%           fig_inter2_combo_counts.png       (bars by combo type)
%           fig_inter3_struct_entropy.png     (N_in/N_out/dH per org)
%           fig_inter4_transition_graph.png   (Hasse layout, one panel per combo)
%  Intra-org (outputs/intra_org_esmo/):
%           fig1_esmo_distributions.png       (violin plots)
%           fig2_feasibility.png              (mean N bars + relative freq)
%           fig3_pc_vs_cp_scatter.png         (PC vs CP scatter)
% ============================================================

\section{Chemical Organization Theory and the Statistical Geometry of Regimes}
\label{sec:cot_lp}

% ------------------------------------------------------------------
\subsection{Organizations, Lattices, and the Conflict Model}
\label{sec:cot_orgs}
% ------------------------------------------------------------------

The dynamical analyses of Sections~\ref{sec:bandtribe}--\ref{sec:chiefdom}
establish what a conflict system \emph{does} under specific parameterisations:
which trajectories emerge, which modes dominate, and at what timescale
strategies become distinguishable.  A complementary and deeper question is
what the system is \emph{structurally capable of}, independent of any
particular choice of rate constants.  Chemical Organization Theory
(COT)~\cite{dittrich2007chemical,veloz2014reaction} provides the language
for this question.

A \textbf{reaction network} consists of a set of species $\mathcal{S}$ and a
set of reactions $\mathcal{R}$, each mapping a multiset of reactants to a
multiset of products, encoded in the stoichiometric matrix
$\mathbf{S} \in \mathbb{R}^{|\mathcal{S}| \times |\mathcal{R}|}$.
A subset $\mathcal{O} \subseteq \mathcal{S}$ is a \textbf{chemical
organisation} if it satisfies two structural conditions simultaneously:
(i)~\emph{closure}---no reaction whose reactants are all in $\mathcal{O}$
produces a species outside $\mathcal{O}$; and
(ii)~\emph{self-maintenance}---there exists a strictly positive flux vector
$\mathbf{v} > \mathbf{0}$ over the reactions enabled by $\mathcal{O}$ such
that the net production of every species in $\mathcal{O}$ is non-negative:
$\mathbf{S}_{\mathcal{O}} \mathbf{v} \geq \mathbf{0}$.
Organizations are thus the subsets of the network that can in principle
sustain themselves indefinitely.  They form a lattice under set inclusion,
visualised as a Hasse diagram in which each edge is a direct containment
step with no intermediate organization.

The conflict model analysed here comprises $|\mathcal{S}| = 14$ species and
$|\mathcal{R}| = 37$ reactions.  Applying COT to this full network yields
23~organizations (nodes $X_1$--$X_{23}$ in
Figure~\ref{fig:cot_hierarchy}), forming a lattice that spans from the
minimal organization $X_1 = \{\text{Res}\}$ at the base to the maximal
organization $X_{23}$ encompassing all species including both foreign-aid
catalysts.  Four of these organizations correspond directly to the
social-structural regimes studied throughout this paper and anchor the
remainder of this section:

\begin{description}
  \item[Tribe ($X_3$)] $\mathcal{T} = \{C,\, C_{\text{Res}},\, G,\,
    \text{Res},\, X\}$ --- the acephalous community with no hierarchical
    species.
  \item[Chief ($X_9$)] $\mathcal{H} = \mathcal{T} \cup \{H,\,
    H_{\text{Res}},\, P_H\}$ --- the chiefdom layer added to the tribal
    base.
  \item[State ($X_{12}$)] $\mathcal{I} = \mathcal{T} \cup \{I,\,
    I_{\text{Res}},\, L,\, P_I\}$ --- the institutional-state layer.
  \item[ChiefState ($X_{20}$)] $\mathcal{CS} = \mathcal{H} \cup
    \mathcal{I} \setminus \mathcal{T}$ --- the dual-hierarchy organization
    containing all twelve conflict-relevant species.
\end{description}

The cover edges of the sub-lattice spanned by these four organizations are
$\mathcal{T} \subset \mathcal{H}$, $\mathcal{T} \subset \mathcal{I}$,
$\mathcal{H} \subset \mathcal{CS}$, and $\mathcal{I} \subset \mathcal{CS}$.
Crucially, neither $\mathcal{H} \subseteq \mathcal{I}$ nor
$\mathcal{I} \subseteq \mathcal{H}$, so no edge exists between Chief and
State in the Hasse diagram --- they are structurally incomparable
organizational forms.

\begin{figure}[ht]
  \centering
  \includegraphics[width=0.80\linewidth]{figures/hierarchy_complex_orgs}
  \caption{Hasse diagram of the 23~chemical organizations of the full
    conflict reaction network (14 species, 37 reactions).
    \textbf{Green nodes} are \emph{self-maintaining} organizations: they
    satisfy both closure and the existence of a strictly positive
    self-sustaining flux.
    \textbf{Cyan nodes} are organizations whose self-maintenance is
    boundary-feasible: they satisfy closure and non-negative net production
    exists, but only at flux configurations where at least one reaction
    operates at its lower bound.
    Arrows denote direct set inclusion (cover edges of the lattice).
    The four regimes used throughout this paper are $X_3$~(Tribe),
    $X_9$~(Chief), $X_{12}$~(State), and $X_{20}$~(ChiefState);
    $X_9$ is the only regime that is strictly self-maintaining without
    external catalytic inputs ($F_H$, $F_I$).}
  \label{fig:cot_hierarchy}
\end{figure}


% ------------------------------------------------------------------
\subsection{Elementary Self-Maintaining Modes of Operation}
\label{sec:esmo_theory}
% ------------------------------------------------------------------

Organizations characterise the \emph{stable substructures} of a reaction
network, but they say nothing about how the system moves \emph{between}
them.  To analyse transitions we introduce the notion of an
\textbf{Elementary Self-Maintaining Mode of Operation} (ESMO).

\subsubsection*{Flux polytopes and process cones}

Let the $O_7$ sub-network denote the 12-species, 35-reaction sub-system
supported by the ChiefState organization $\mathcal{CS}$ --- the common
ambient space for all four regimes.  Its stoichiometric matrix is
$\mathbf{S} \in \mathbb{R}^{12 \times 35}$.  Admissible flux states are
non-negative vectors $\mathbf{v} \in \mathbb{R}^{35}_{\geq 0}$,
i.e.\ reactions can only fire forward.  This defines the \emph{positive
orthant cone} $\mathcal{K} = \mathbb{R}^{35}_{\geq 0}$.

For a directed transition from an organizational regime $A$ to a regime $B$,
we define a \textbf{transition signature}
$\boldsymbol{\sigma}_{A \to B} \in \{-1,\, 0,\, +1\}^{12}$ derived from
the characteristic states of the two regimes.  Let $\mathbf{x}^A$ and
$\mathbf{x}^B$ be representative species-abundance vectors for regimes $A$
and $B$ respectively.  For each species $s$:
\begin{equation}
  \sigma_s = \begin{cases}
    +1  & \text{if } x^B_s - x^A_s >  \theta \\
    -1  & \text{if } x^B_s - x^A_s < -\theta \\
    \phantom{+}0 & \text{otherwise,}
  \end{cases}
  \label{eq:signature}
\end{equation}
where $\theta > 0$ is a neutrality threshold that filters species whose
abundance difference is negligible relative to the regime contrast.
The signature encodes a \emph{coverability constraint} in the sense of
reachability theory~\cite{veloz2022resilience}: species with
$\sigma_s = +1$ must be produced in net excess ($[\mathbf{S}\mathbf{v}]_s > 0$),
and species with $\sigma_s = -1$ must be consumed in net excess
($[\mathbf{S}\mathbf{v}]_s < 0$).  This is a \emph{directional} condition,
not a quantitative one: it asks whether the stoichiometry supports movement
in the right direction, not whether it reaches the exact target state.

Because every reaction in the network converts one person-type to another
without creating or destroying people, the total population
$P = C + X + H + I + P_H + P_I$ is a conserved quantity.  This structural
constraint is encoded in the left null-vector of the stoichiometric matrix
and must be respected when constructing characteristic states for cross-scale
transitions.  All representative states used in the inter-organizational
analysis are therefore normalized to a common reference population
$P_{\text{ref}} = 1000$ before computing signatures, ensuring that the
LP constraints are satisfiable by construction.

The constrained cone for transition $A \to B$ is
\begin{equation}
  \mathcal{C}_{A \to B} \;=\; \bigl\{ \mathbf{v} \in \mathbb{R}^{35}_{\geq 0}
  \;\big|\; \sigma_s [\mathbf{S}\mathbf{v}]_s > 0 \;\; \forall\, s:\, \sigma_s \neq 0
  \bigr\},
\end{equation}
intersected with the simplex $\mathbf{1}^{\top}\mathbf{v} = 1$ to obtain
a bounded feasible polytope $\mathcal{P}_{A \to B}$.  An ESMO is a
\textbf{vertex} of $\mathcal{P}_{A \to B}$: a minimal, non-decomposable
reaction coalition that simultaneously satisfies all sign constraints.
Vertices of a linear polytope are the elementary solutions from which all
feasible flux states can be written as convex combinations; they are the
direct structural analogue of elementary flux modes in metabolic
stoichiometry~\cite{schuster1994elementary}.

\subsubsection*{ESMO count as structural entropy}

The number of vertices $N(A \to B) = |\text{vert}(\mathcal{P}_{A \to B})|$
measures the \emph{microstate count} for the transition macrostate defined
by signature $\boldsymbol{\sigma}_{A \to B}$.  In the spirit of Boltzmann,
a larger $N$ means more distinct minimal reaction coalitions can realize the
required directional change: the transition is structurally richer.
Conversely, if $\mathcal{P}_{A \to B} = \emptyset$ the transition is
stoichiometrically infeasible regardless of kinetics.

This is not a statement about probability in the kinetic sense --- a
transition with $N = 2000$ ESMOs may still be dynamically rare if all
those ESMOs require co-activation of reactions with small rate constants.
What it provides is a \textbf{parameter-free lower bound on structural
accessibility}: a transition with $N = 0$ (infeasible polytope) is
impossible for any parameterisation; a transition with large $N$ has many
distinct structural pathways and is in that sense \emph{entropically
favoured} at the level of network topology.

\subsubsection*{Non-parametric probability and the Markovian framework}

When multiple representative states $\mathbf{x}^A_1, \ldots,
\mathbf{x}^A_m$ and $\mathbf{x}^B_1, \ldots, \mathbf{x}^B_n$ are
available for each regime, the pairwise ESMO computations produce a
\emph{distribution} of vertex counts over all $m \times n$ signature
problems.  Two quantities of interest emerge without any kinetic
assumptions:
\begin{equation}
  \Pr(A \to B) \;=\; \frac{|\{(i,j)\,:\, N(\mathbf{x}^A_i \to \mathbf{x}^B_j) > 0\}|}{m \cdot n},
  \quad
  \bar{N}(A \to B) \;=\; \frac{1}{mn}\sum_{i,j} N(\mathbf{x}^A_i \to \mathbf{x}^B_j).
  \label{eq:esmo_prob}
\end{equation}
$\Pr(A \to B)$ is the \textbf{structural feasibility probability}: the
fraction of state-pair microstates for which the transition is realizable
at all.  It depends only on whether the polytope is non-empty, making it
the most conservative and robust claim.  $\bar{N}(A \to B)$ is the
\textbf{mean structural accessibility}: the expected number of minimal
pathway coalitions available, a richer but still fully parameter-free
quantity.  If $\Pr(A \to B) > \Pr(B \to A)$ we say $A \to B$ is
\emph{structurally dominant}; if additionally $\bar{N}(A \to B) \gg
\bar{N}(B \to A)$ the dominance is strong.

This connects directly to the Markovian framework developed
in~\citet{veloz2023markovian}, where the organizations of a reaction
network are treated as nodes of a stochastic process and transition
probabilities encode the tendency of the system to move from one
organizational attractor to another.  The ESMO distribution provides a
principled, data-observable basis for estimating these transition
probabilities without fitting kinetic parameters: states are far easier to
measure empirically (population counts, resource levels, institutional
indicators) than reaction rate constants.
The structural entropy index of an organization,
$\Delta H_{\mathcal{O}} = \log \overline{N}_{\text{in}}(\mathcal{O}) -
\log \overline{N}_{\text{out}}(\mathcal{O})$,
where $\overline{N}_{\text{in}}$ and $\overline{N}_{\text{out}}$ aggregate
mean accessibility over all incoming and outgoing cover-edge transitions
respectively, provides a single scalar measure of whether the organization
acts as a structural attractor ($\Delta H > 0$, centripetal) or as a
structural transient ($\Delta H < 0$, centrifugal) in the Markovian process
on the organizational lattice.


% ------------------------------------------------------------------
\subsection{Intra-Organizational Transition Probabilities}
\label{sec:esmo_intra}
% ------------------------------------------------------------------

The inter-organizational analysis of Section~\ref{sec:esmo_inter} uses a
single characteristic state per organizational type.  This is sufficient for
establishing directional asymmetry, but it raises a legitimate concern: do
the results depend on the particular choice of representative state?  The
multi-state ensemble analysis addresses this directly.

\subsubsection*{Multi-state ensemble design}

For each organizational regime and each qualitative type (peace / conflict),
we define a library of $m = n = 5$ representative states, carefully designed
so that different pairs $(i, j)$ produce distinct LP sign patterns (signatures).
The \textbf{non-proportionality criterion} requires that states differ in
\emph{which} species cross the neutrality threshold $\theta$, not merely in
overall abundance level: proportional rescaling of a state would yield the
same LP signature and add no statistical information.  For example, among the
five Tribe peace states, one has an elevated warrior species ($X = 120$) that
falls \emph{below} the conflict levels for some conflict states, generating
a unique signature in which $X$ must \emph{shrink} on the transition to
conflict --- the opposite of the typical direction.

All $m \times n = 25$ PC pairs and all $n \times m = 25$ CP pairs are
enumerated per organisation, yielding 200 directed LP problems in total.
Identical LP signatures are detected by a hash of the sign vector and solved
only once (cache deduplication), so the computational cost scales with the
number of distinct signatures, not the number of pairs.

From the distribution of ESMO counts over all pairs, we report:
\begin{align*}
  \Pr_{\text{PC}}(\mathcal{O}) &= \frac{|\{(i,j): N(\mathbf{x}^P_i \to \mathbf{x}^C_j) > 0\}|}{25}, &
  \bar{N}_{\text{PC}}(\mathcal{O}) &= \frac{1}{25}\sum_{i,j} N(\mathbf{x}^P_i \to \mathbf{x}^C_j),
\end{align*}
and symmetrically for CP.  Both quantities are fully parameter-free:
$\Pr_{\text{PC}}$ depends only on polytope non-emptiness,
while $\bar{N}_{\text{PC}}$ is a structural richness estimate.
Because $\Pr_{\text{PC}} = \Pr_{\text{CP}} = 1.0$ for all four organizations
(every state-pair combination yields a feasible polytope), the informative
comparison across organizations is the \emph{relative frequency} of each
direction:
\begin{equation}
  f_{\text{PC}}(\mathcal{O}) = \frac{\bar{N}_{\text{PC}}}{\bar{N}_{\text{PC}} + \bar{N}_{\text{CP}}},
  \qquad
  f_{\text{CP}}(\mathcal{O}) = \frac{\bar{N}_{\text{CP}}}{\bar{N}_{\text{PC}} + \bar{N}_{\text{CP}}},
  \label{eq:rel_freq}
\end{equation}
which measures the structural weight of each transition direction independently
of absolute ESMO counts.

\subsubsection*{Tribe: a complete example}

Table~\ref{tab:tribe_intra} reports the Tribe results over all 50 pairs.
Both peace-to-conflict and conflict-to-peace transitions are \emph{fully
feasible}: $\Pr_{\text{PC}} = \Pr_{\text{CP}} = 1.0$, meaning that for
every combination of representative states the stoichiometry supports
the directional change.  This confirms that the Tribe sub-network is
topologically capable of cycling between peace and conflict regardless of
the specific state realization, consistent with its role as a baseline
organizational form.

The mean ESMO counts reveal a directional bias:
$\bar{N}_{\text{PC}} \approx 1140 > \bar{N}_{\text{CP}} \approx 899$,
a ratio of approximately $1.27$, corresponding to relative frequencies
$f_{\text{PC}} = 0.56$ and $f_{\text{CP}} = 0.44$.
Across the 25 PC pairs, six distinct LP signatures are realized
(range $711$--$1251$), and likewise six for CP (range $740$--$1114$),
confirming that the ensemble genuinely samples distinct structural
microstates rather than collapsing to a single polytope.
The standard deviations ($\sigma_{\text{PC}} \approx 185$,
$\sigma_{\text{CP}} \approx 90$) quantify the spread due to the choice of
representative state.  Their modest size relative to the means (16\% and
10\% coefficients of variation respectively) indicates that the direction of
the inequality $\bar{N}_{\text{PC}} > \bar{N}_{\text{CP}}$ is robust to
state selection.

\begin{table}[ht]
  \centering
  \caption{Multi-state ESMO ensemble statistics for Tribe (5 peace
    $\times$ 5 conflict states; 25 PC and 25 CP directed pairs).}
  \label{tab:tribe_intra}
  \begin{tabular}{lrrrrcr}
    \hline
    Direction & $\Pr_{\text{feasible}}$ & $\bar{N}$ & $\sigma_N$ &
    Distinct sigs & Range & $f$ \\
    \hline
    PC (peace $\to$ conflict) & 1.000 & 1139.8 & 184.6 & 6 & 711--1251 & 0.559 \\
    CP (conflict $\to$ peace) & 1.000 &  899.0 &  89.5 & 6 & 740--1114 & 0.441 \\
    \hline
  \end{tabular}
\end{table}

Structurally, the Tribe community is easier to destabilize than to
recover: the parameter-free ratio $\bar{N}_{\text{PC}} /
\bar{N}_{\text{CP}} \approx 1.27$ holds robustly across the state
ensemble.  This complements the simulation results of
Section~\ref{sec:bandtribe}, where conflict escalation required fewer
interventions than recovery under the same resource conditions.

\subsubsection*{Chief, State, and ChiefState}

The pattern observed in Tribe extends to all four organizational regimes,
but with a striking amplification as complexity grows.
Table~\ref{tab:all_intra} reports the full ensemble statistics for all four
organizations.

\begin{table}[ht]
  \centering
  \caption{Multi-state ESMO ensemble statistics for all four organizational
    regimes (5 peace $\times$ 5 conflict states; 25 PC and 25 CP pairs per
    organization; all $\Pr_{\text{feasible}} = 1.000$).
    $f$ denotes the relative frequency $\bar{N}_{\text{combo}} /
    (\bar{N}_{\text{PC}} + \bar{N}_{\text{CP}})$.}
  \label{tab:all_intra}
  \begin{tabular}{llrrrrr}
    \hline
    Organization & Dir. & $\bar{N}$ & $\sigma_N$ & Distinct sigs & Range & $f$ \\
    \hline
    Tribe & PC & 1139.8 & 184.6 & 6  & 711--1251   & 0.559 \\
          & CP &  899.0 &  89.5 & 6  & 740--1114   & 0.441 \\
    \hline
    Chief & PC & 1502.7 & 315.4 & 11 & 883--1843   & 0.596 \\
          & CP & 1018.8 & 139.2 & 11 & 739--1258   & 0.404 \\
    \hline
    State & PC & 1606.4 & 191.1 & 14 & 1145--1953  & 0.698 \\
          & CP &  696.6 & 146.0 & 14 & 380--913    & 0.302 \\
    \hline
    ChiefState & PC & 1804.6 & 215.8 & 22 & 1307--2217 & 0.678 \\
               & CP &  858.7 & 152.2 & 22 & 526--1118  & 0.322 \\
    \hline
  \end{tabular}
\end{table}

Several features of Table~\ref{tab:all_intra} warrant discussion.

\paragraph{The PC/CP ratio gradient.}
The ratio $\bar{N}_{\text{PC}} / \bar{N}_{\text{CP}}$ increases
monotonically with organizational complexity: $1.27$ (Tribe), $1.47$
(Chief), $2.10$ (ChiefState), and $2.31$ (State).  In terms of relative
frequencies, the structural weight of destabilization grows from 56\%
(Tribe) to 60\% (Chief) to 70\% (State) and 68\% (ChiefState).
This is a parameter-free topological result: more complex organizational
forms have a progressively stronger structural bias toward transitioning
from peace to conflict than in the reverse direction.  No rate constant
calibration is required to observe this gradient; it is an intrinsic
property of the conflict network stoichiometry at each organizational scale.

\paragraph{Signature diversity and robustness.}
The number of distinct LP signatures realised per direction also scales with
organizational complexity: 6 for Tribe, 11 for Chief, 14 for State, and 22
for ChiefState.  A higher signature count means the state ensemble samples
genuinely distinct structural microstates, not just quantitative variants of
the same qualitative constraint.  The coefficients of variation
($\sigma / \bar{N}$) remain moderate in all cases --- at most 21\%
for Chief PC --- confirming that the directional inequalities are robust
to state selection even when the distributions have non-trivial spread.

\paragraph{State and ChiefState extremes.}
The State organization exhibits the sharpest intra-organizational asymmetry
($\bar{N}_{\text{PC}} / \bar{N}_{\text{CP}} = 2.31$): peace-to-conflict
transitions average 1606 ESMOs while conflict-to-peace transitions average
only 697.  ChiefState shows a similar but slightly smaller asymmetry (ratio
2.10), with the highest absolute mean ESMO counts in both directions
($\bar{N}_{\text{PC}} = 1805$, $\bar{N}_{\text{CP}} = 859$), reflecting
the richer reaction set available in the dual-hierarchy organization.

\paragraph{Figures.}
Figure~\ref{fig:intra_distributions} displays the full ESMO count
distributions for each organization as violin plots on a shared axis,
making the within-organization spread and the between-direction asymmetry
directly comparable.  Figure~\ref{fig:intra_feasibility} compares the
absolute mean ESMO counts across organizations and directions, annotated
with their relative frequencies.  Figure~\ref{fig:intra_scatter} shows
matched PC-vs-CP scatter plots: each point represents one state pair $(i,j)$
with its N value in both directions; points above the diagonal correspond to
pairs where recovery (CP) is structurally richer than destabilisation (PC),
and point density (encoded by overlap transparency) reveals that such
reversals are rare in all four organizations.

\begin{figure}[ht]
  \centering
  \includegraphics[width=\linewidth]{figures/fig1_esmo_distributions}
  \caption{Distribution of ESMO counts across all 25 state pairs per
    organization and direction (intra-organizational transitions).
    Violin plots share a common $y$-axis; the white dot marks the mean.
    PC (orange) = peace$\to$conflict (destabilisation);
    CP (blue) = conflict$\to$peace (recovery).
    All 200 polytopes are feasible ($\Pr = 1.0$).
    The progressive upward shift of the PC violin and downward shift of
    the CP violin as organizational complexity increases reflects the
    growing structural bias toward destabilisation.}
  \label{fig:intra_distributions}
\end{figure}

\begin{figure}[ht]
  \centering
  \includegraphics[width=0.85\linewidth]{figures/fig2_feasibility}
  \caption{Mean ESMO count per direction and organization (bars $\pm$ 1 SD).
    Bar annotations give the absolute mean and the relative frequency
    $f = \bar{N} / (\bar{N}_{\text{PC}} + \bar{N}_{\text{CP}})$.
    The right axis shows the relative frequency scale.
    The systematic dominance of PC over CP, and its amplification with
    organizational complexity, is evident across all four regimes.}
  \label{fig:intra_feasibility}
\end{figure}

\begin{figure}[ht]
  \centering
  \includegraphics[width=\linewidth]{figures/fig3_pc_vs_cp_scatter}
  \caption{Matched PC vs.\ CP ESMO counts per state pair (intra-organizational).
    Each point corresponds to one $(i,j)$ pair; its $x$-coordinate is
    $N_{\text{PC}}$ and its $y$-coordinate is $N_{\text{CP}}$.
    The dashed diagonal marks $N_{\text{PC}} = N_{\text{CP}}$.
    Points above the diagonal have more recovery pathways than
    destabilisation pathways; points below the diagonal have more
    destabilisation pathways.
    Overlap transparency encodes point density: darker regions indicate
    many pairs sharing the same $(N_{\text{PC}}, N_{\text{CP}})$ value.
    In all four organizations the majority of pairs fall below the diagonal,
    confirming the structural dominance of destabilisation.}
  \label{fig:intra_scatter}
\end{figure}


% ------------------------------------------------------------------
\subsection{Inter-Organizational Transitions along the Hasse Diagram}
\label{sec:esmo_inter}
% ------------------------------------------------------------------

We computed ESMOs for every directed pair along the four cover edges of the
sub-lattice $\{\mathcal{T}, \mathcal{H}, \mathcal{I}, \mathcal{CS}\}$,
in both directions and across all four peace/conflict state combinations
(PP, PC, CP, CC), yielding 32 structurally distinct LP problems in the
12-dimensional flux space of the $O_7$ sub-network.  The complete ESMO
count matrices for each directed edge are shown in
Figure~\ref{fig:esmo_matrices}; the aggregated counts by transition type
across all edges are summarised in Figure~\ref{fig:esmo_combos}.  A directed
graph summary of the lattice, with edge thickness proportional to the square
root of the ESMO count for each combo, is given in
Figure~\ref{fig:esmo_graph}.

\begin{figure}[ht]
  \centering
  \includegraphics[width=\linewidth]{figures/fig_inter4_transition_graph}
  \caption{Directed ESMO transition graph for the four organizational
    regimes, one panel per state-combination type (PP, PC, CP, CC).
    Nodes are positioned in a diamond Hasse layout: Tribe at the base,
    Chief and State at mid-level, ChiefState at the apex.  Arrow curvature
    distinguishes upward (arc right) from downward (arc left) transitions
    along each cover edge.  Arrow width scales as $\sqrt{N}$, where $N$ is
    the ESMO count; the label reports the exact $N$ value.  Width
    asymmetry between opposite arrows along the same edge reveals the
    directional structural bias discussed in the text.}
  \label{fig:esmo_graph}
\end{figure}

\begin{figure}[ht]
  \centering
  \includegraphics[width=\linewidth]{figures/fig_inter1_heatmaps}
  \caption{ESMO count matrices for the four cover edges of the
    conflict organizational lattice.  Each $2\times 2$ cell shows the
    number of ESMO vertices for a directed transition, with rows
    indicating the source state (peace / conflict) and columns the
    target state.  Left column: going \emph{up} the lattice (smaller
    $\to$ larger organization); right column: going \emph{down}.
    Colour intensity encodes ESMO count on a common scale across all
    eight panels (yellow $\approx 380$, dark red $\approx 2153$).}
  \label{fig:esmo_matrices}
\end{figure}

\begin{figure}[ht]
  \centering
  \includegraphics[width=\linewidth]{figures/fig_inter2_combo_counts}
  \caption{ESMO counts disaggregated by state-combination type
    (PP = peace$\to$peace, PC = peace$\to$conflict,
    CP = conflict$\to$peace, CC = conflict$\to$conflict)
    for each cover edge.  Left panel: going up the lattice ($\to$ larger
    organization); right panel: going down.  Bar height encodes the ESMO
    count; taller bars indicate greater structural richness.  The PC bars
    consistently dominate the going-down panels, and CC bars dominate
    the going-up panels for the ChiefState edges.}
  \label{fig:esmo_combos}
\end{figure}

\subsubsection*{Structural entropy and the attractor landscape}

Computing $\Delta H_{\mathcal{O}}$ for each of the four organizations from
the 32 directed ESMO counts reveals a clear hierarchy of structural
attraction (Table~\ref{tab:struct_entropy};
Figure~\ref{fig:esmo_struct_entropy}).

Tribe is the only centripetal organization ($\Delta H_{\mathcal{T}} =
+0.414$): its high mean incoming accessibility
($\overline{N}_{\text{in}} = 1261$, the largest among all four orgs)
reflects the abundance of minimal reaction coalitions that drive complex
organizations downward into the tribal base, while its low outgoing
accessibility ($\overline{N}_{\text{out}} = 834$, the smallest) reflects
the stoichiometric difficulty of building institutional layers from a tribal
starting point.  Tribe is thus the unique \emph{structural attractor} in
the Markovian process on the lattice, consistent with its role as the robust
baseline regime in the simulations of
Sections~\ref{sec:bandtribe}--\ref{sec:chiefdom}.

Chief is very weakly centrifugal ($\Delta H_{\mathcal{H}} = -0.048$),
essentially balanced: its incoming and outgoing mean accessibilities differ
by less than 5\% ($\overline{N}_{\text{in}} = 1146$,
$\overline{N}_{\text{out}} = 1202$).  Chief is a high-throughput conduit
--- it receives and emits a large number of structural pathways --- but
without a systematic directional preference.  This structural neutrality
may reflect the intermediate position of the chiefdom in the lattice:
equally connected upward (to ChiefState) and downward (to Tribe), it
acts as a transition hub rather than an attractor.

ChiefState is centrifugal ($\Delta H_{\mathcal{CS}} = -0.136$): despite
the highest outgoing accessibility ($\overline{N}_{\text{out}} = 1234$,
driven by the globally maximal $N = 2153$ for the ChiefState-to-State
descent), its incoming accessibility is lower
($\overline{N}_{\text{in}} = 1077$), reflecting the stoichiometric
constraints on assembling dual-hierarchy structures from below.
ChiefState is a structural transient: the network topology does not
favour its formation but generously supports its transformation into
other organizational forms.

The most clear-cut centrifugal organization is State
($\Delta H_{\mathcal{I}} = -0.209$), which has the lowest mean incoming
accessibility ($\overline{N}_{\text{in}} = 922$) --- dominated by the
near-infeasible Tribe-to-State upward pathways (minimum $N = 380$ for
the CC combo, $N = 521$--$631$ for the others) --- while its outgoing
accessibility ($\overline{N}_{\text{out}} = 1137$) benefits from the
large ESMO counts on State-to-Tribe descents ($N = 966$--$1666$).
State is a structural transient: the network topology does not favour its
emergence but supports its dissolution in multiple directions.

A striking feature of the normalized results is that \emph{all} organizations
except Tribe are centrifugal ($\Delta H < 0$).  Under population-consistent
normalization, the organizational lattice is fundamentally asymmetric: only
the minimal Tribe regime acts as a structural attractor, while Chief,
ChiefState, and State all exhibit net structural outflow.  This is consistent
with the evolutionary pressure toward tribal-level organization observed in the
simulation collapse scenarios of Section~\ref{sec:chiefdom}.

\begin{table}[ht]
  \centering
  \caption{Structural entropy indices for the four organizational regimes,
    computed from the 32 directed ESMO counts across cover edges.
    $\overline{N}_{\text{in}}$: mean ESMO count over all incoming transitions;
    $\overline{N}_{\text{out}}$: mean ESMO count over all outgoing transitions;
    $\Delta H = \ln\overline{N}_{\text{in}} - \ln\overline{N}_{\text{out}}$:
    positive = centripetal (structural attractor), negative = centrifugal
    (structural transient).}
  \label{tab:struct_entropy}
  \begin{tabular}{lrrr}
    \hline
    Organization   & $\overline{N}_{\text{in}}$ & $\overline{N}_{\text{out}}$ & $\Delta H$ \\
    \hline
    Tribe          & 1261.2 &  834.0 & $+0.414$ \\
    Chief          & 1145.6 & 1201.5 & $-0.048$ \\
    ChiefState     & 1077.0 & 1233.8 & $-0.136$ \\
    State          &  922.1 & 1136.8 & $-0.209$ \\
    \hline
  \end{tabular}
\end{table}

\begin{figure}[ht]
  \centering
  \includegraphics[width=0.88\linewidth]{figures/fig_inter3_struct_entropy}
  \caption{\textbf{Top:} Mean incoming ($\overline{N}_{\text{in}}$, blue) and
    outgoing ($\overline{N}_{\text{out}}$, orange) ESMO counts for each
    organizational regime, averaged over all 8 directed transitions entering
    or leaving that regime.  Numbers above each bar give the exact values.
    \textbf{Bottom:} Structural entropy index
    $\Delta H_{\mathcal{O}} = \ln\overline{N}_{\text{in}} -
    \ln\overline{N}_{\text{out}}$ for each regime.
    Green bar ($\Delta H > 0$): centripetal organizational attractor, with
    more incoming than outgoing structural pathways.
    Red bars ($\Delta H < 0$): centrifugal transients, with more outgoing
    than incoming pathways.
    Tribe is the sole centripetal regime; Chief, ChiefState, and State
    are all centrifugal under population-consistent normalization.}
  \label{fig:esmo_struct_entropy}
\end{figure}

These results complement the mode-projection analysis of
Section~\ref{sec:modes} by providing a parameter-free structural
explanation for the dynamical patterns observed there.  The exploitation
trajectory's tendency to collapse toward tribal-level organization is not
merely a consequence of the chosen rate constants: the organizational
lattice structurally favours downward transitions under conflict conditions,
particularly along the State and ChiefState edges.  This structural bias
is a topological property of the conflict reaction network, visible only
through the ESMO lens.
